% v1_draft.tex — Ceremonial Entheogen Use in Puerto Rico
% Phase 3 drafting output. Built on manuscript/drafts/v0_jpd_template.tex.
% Every number traces to outputs/report_v2.md (N=73 analytic sample).
% Formatting per manuscript/JOURNAL_REQUIREMENTS.md (Journal of Psychoactive Drugs).
% Word budget: Introduction-Methods-Results-Discussion <= 4,000 words (strict, per author decision).

\documentclass[12pt]{article}

\usepackage[utf8]{inputenc}
\usepackage[T1]{fontenc}
\usepackage{times}
\usepackage[margin=1in]{geometry}
\usepackage{setspace}
\doublespacing
\usepackage{lineno}
\usepackage{graphicx}
\usepackage{booktabs}
\usepackage{amsmath}
\usepackage{url}
\urlstyle{same}  % render URLs in the body font; avoids a Courier dependency
\emergencystretch=3em  % double spacing leaves little slack; prevents overfull lines at long citations
\usepackage[authoryear,round]{natbib}
\bibliographystyle{chicago}

\title{Ceremonial Entheogen Use in Puerto Rico: A Descriptive Survey of Participants and Facilitators}

\author{
  Yamil O. Ortiz Ortiz$^{1}$ \and Jean C. V\'elez Rodr\'iguez$^{2}$ \and Juli\'an M. Hern\'andez Torres$^{2,3}$ \and
  Juliana Mill\'an-Torres$^{4}$ \and Paulina Rull\'an Farinacci$^{5}$ \and Carolina Riollano$^{6}$ \and
  Ang\'elica N. Grajales Rold\'an$^{1}$ \and Adriana I. Rodr\'iguez Massa$^{1}$ \\[1em]
  \begin{minipage}{0.92\textwidth}\centering\small\setstretch{1.15}
  $^{1}$Centro de Investigaciones Sociales, Facultad de Ciencias Sociales, University of Puerto Rico,
  R\'io Piedras Campus, San Juan, Puerto Rico \\
  $^{2}$Puerto Rico Institute for Psychedelic Science, Medicine \& Awareness, San Juan, Puerto Rico \\
  $^{3}$Behavioral Sciences Research Institute, University of Puerto Rico, Medical Sciences Campus,
  San Juan, Puerto Rico \\
  $^{4}$Colectivo Psicod\'elico de Puerto Rico, Colectivo Psicod\'elico Inc., and Within Psychological
  Services, San Juan, Puerto Rico \\
  $^{5}$Department of Psychiatry, School of Medicine, University of Puerto Rico, Medical Sciences Campus,
  San Juan, Puerto Rico \\
  $^{6}$Colectivo Psicod\'elico de Puerto Rico, San Juan, Puerto Rico \\[0.5em]
  Corresponding author: Jean C. V\'elez Rodr\'iguez, jean.velez@prismapr.org
  \end{minipage}
}

\date{}

\begin{document}

\maketitle
\linenumbers

\begin{abstract}
\noindent
% ABSTRACT — written last, per plan. Unstructured, <=200 words (JPD requirement).
% Draft below is provisional and will be finalized after Phase 5 review.
Ceremonial use of entheogens occurs in Puerto Rico outside clinical and regulatory structures, and the
practices surrounding it have not been described. We conducted an anonymous, bilingual, cross-sectional
online survey of adults aged 21 and over reporting involvement in entheogen-assisted spiritual ceremonies
in Puerto Rico, recruited by convenience and chain referral between February and June 2026. Of 100
submissions, 73 respondents identified a role: 72 as ceremony participants, 8 as facilitators, and 7 as
both. Respondents were middle-aged (mean 42.4 years), predominantly university-educated, and almost all
resident in Puerto Rico. Most recent ceremonies were typically single, multi-hour events in natural
settings involving psilocybin mushrooms or ayahuasca, and satisfaction was strongly ceiling-weighted
(median 6 of 6). Preparation was near-universal among respondents answering, while pre-ceremony screening
was not, and the small facilitator subsample described training and emergency planning as informal rather
than credentialed. Five of 67 participants answering reported non-consensual physical contact during a
ceremony. A pre-declared comparison of any preparation against none on satisfaction was directionally
positive but rests on five unprepared respondents and is exploratory. These findings describe this sample
only and support no prevalence or causal claim.
\end{abstract}

\noindent\textbf{Keywords:} entheogens; psychedelics; ceremonial use; Puerto Rico; harm reduction; survey research

\section{Introduction}

Ceremonial use of entheogens---psilocybin-containing mushrooms, ayahuasca, and related preparations taken
in structured group or guided settings---has become substantially more visible across the Americas over the
past two decades, largely outside the clinical-trial infrastructure that generates most of the contemporary
research literature \citep{golden2022effects,carvalho2025scoping}. The resulting evidence base is
asymmetric. Randomized clinical research characterizes carefully screened participants receiving
standardized doses under monitored conditions, and its samples remain narrow
\citep{hughes2024ethnoracial,marcus2026psychedelics}, while community ceremonial practice is documented
mainly through ethnographic and qualitative work, with comparatively few survey-level descriptions of who
takes part and under what conditions \citep{carvalho2025scoping}.

Puerto Rico occupies an unusual position in this landscape. It is a Spanish-speaking Caribbean jurisdiction
in which psilocybin, DMT, and related compounds are Schedule I substances under the federal Controlled
Substances Act \citep{uscongress1970csa} and under the territory's own aligned statute \citep{pr1971csa},
and which sits outside the state-level decriminalization and regulated-access reforms that have reshaped
the legal environment in parts of the United States \citep{siegel2023psychedelic}. It is also a setting
with a long-documented tradition of spiritually framed healing practice---espiritismo and related
practices---that has been described in the clinical and public health literature for more than four decades
\citep{koss1980therapist,comasdiaz1981puertorican,bird1981sociopsychiatry,zerrate2022espiritismo}.

Local survey evidence on entheogen use is nonetheless minimal. One published Puerto Rico--specific survey
exists, conducted by members of the present research team, but it is narrower in scope: it characterizes
psilocybin use patterns, motivations, and personality correlates rather than ceremonial practice, and
addresses neither facilitation nor safety \citep{velezrodriguez2026psilocybin}. An unpublished thesis has
examined university-community perceptions of psilocybin mushrooms \citep{doeltermorales2020hongos}. Neither
describes how ceremonies are actually conducted. Basic descriptive questions---who takes part, which
substances are used, where and for how long ceremonies occur, whether participants are prepared or screened
beforehand, and what adverse experiences are reported---therefore have no published local answer.

A descriptive characterization contributes in three ways. First, it establishes whether practices reported
in a community sample resemble the harm-reduction structures assumed by clinical and retreat-based
frameworks: preparation, pre-ceremony screening, and a plan for difficult experiences
\citep{gorman2021psychedelic,pilecki2021ethical,dutton2025harm}. Second, it identifies which safety
outcomes---including boundary violations, shown elsewhere to be measurable in naturalistic samples
\citep{kruger2025psychedelic,peluso2020reflections}---occur at frequencies justifying measurement in a
larger study. Third, it produces the instrument-level evidence needed before any confirmatory study is
designed in this population.

Our objectives were to describe the sample and the ceremonial context, and to examine bivariate
associations between preparation and reported ceremony outcomes and between participant-observable
facilitator practices and perceived safety. A safety and adverse-event domain, absent from the original
protocol, was added during analysis. Every objective is descriptive. The design supports no statement about
how common these practices are in Puerto Rico, and no statement about what any practice produces.

\section{Methods}

\subsection{Design, ethics, and consent}

This was a single-timepoint, cross-sectional, self-administered online survey of adults reporting
involvement in entheogen-assisted spiritual ceremonies in Puerto Rico. The sample was non-probability and
the analysis unweighted. The study was approved by the University of Puerto Rico Institutional Review Board
(protocol 2509472552, approved 4 November 2025) and was fielded on KoboToolbox between 24 February and 9
June 2026. Participation was fully anonymous: no names, dates of birth, or other personal identifiers were
collected. Respondents read an information sheet covering purpose, procedures, risks, voluntariness, and
data retention, and indicated consent by clicking through to the questionnaire; no signed consent document
was collected, consistent with the anonymous design. ``Prefer not to answer'' was offered on every item,
withdrawal was permitted at any point, no compensation was offered, and mental-health support resources
were listed. An optional future-contact form was linked separately, with no join key to questionnaire
responses. Reporting follows STROBE for cross-sectional studies (Supplementary Material).

\subsection{Recruitment and sample construction}

Recruitment used convenience sampling with chain referral, promoted online through Facebook, Instagram,
Twitter/X, and WhatsApp, together with bulletin-board announcements at the University of Puerto Rico R\'io
Piedras and Medical Sciences campuses. Inclusion criteria were age 21 or older, ability to read and write
in Spanish, access to an internet-connected device, and prior participation in an entheogen-assisted
spiritual ceremony in Puerto Rico. Persons under 21 or unable to consent were excluded. The recruitment
target of 100 responses was met.

Of the 100 submissions collected, 13 were abandoned before any substantive item and were excluded, leaving
87 respondents who reached the role-identification questions. Of those, 14 (16\%) answered ``no'' to both
role items, indicating no ceremony involvement, and are reported here as a screening yield rather than as a
described subgroup. The analytic sample is the remaining $N=73$ respondents who identified a role. Roles
are not a partition: 72 respondents identified as ceremony participants and 8 as facilitators, with 7
identifying as both and appearing in both role-specific analyses. Critically, the two branches are
\emph{unlinked}. The instrument carries no ceremony or facilitator identifier, so no participant can be
connected to the facilitator who ran their ceremony, and facilitator-reported and participant-reported
measures cannot be crossed. Figure~\ref{fig:flow} shows the full construction.

\subsection{Instrument}

No validated Puerto Rico--specific instrument for characterizing ceremonial entheogen use exists, so the
questionnaire was developed by the research team, reviewed iteratively, professionally translated by direct
translation, and given a final team review for cross-language consistency. It was deployed as a single
bilingual XLSForm with Spanish as the default language and English available through the platform's
language switcher, rather than as two separate instruments; open-text responses were accordingly
mixed-language.

The primary outcome, satisfaction with the most recent ceremony, asked ``From 0 to 6, how pleasant was the
ceremony?'' with anchors from \emph{very unpleasant} (0) to \emph{very pleasant} (6), and was treated as
ordinal. The primary exposure asked, ``Aside from the facilitator's instructions or recommendations, do you
personally engage in preparation for these ceremonies?'' (yes / sometimes / no), and therefore captures
self-directed preparation only, not facilitator-provided instruction. Secondary participant-side measures
covered in-ceremony comfort and safety (0--3), trust in the facilitator (0--4), prior knowledge of the
facilitator's background, whether the facilitator asked screening questions beforehand, and whether the
respondent had spoken to a health professional beforehand. The non-consensual contact item
asked whether the respondent had ever experienced physical touch during a ceremony that felt unexpected or
that they had not clearly agreed to beforehand, explicitly including moments of uncertainty or absent
recall; its breadth exceeds a narrower assault framing and the resulting figure should be read
accordingly. Facilitator-branch items covered training, screening approach, distress-response protocol,
emergency planning, and crisis signs witnessed. Three substance variables were kept distinct and never
merged: facilitator-led substances, participant substances across all ceremonies attended, and the
substance used in the most recent ceremony. Attention checks were branch-specific. Verbatim wording and
anchors for all key items appear in the Supplementary Material.

\subsection{Missing data, structural skips, and suppression}

Denominators shift across items for two distinct reasons, kept separate throughout: structural skips arise
from branch logic and are reported as their own row, while item nonresponse is the residual. Every figure
carries its own base and no denominator is carried across sections. ``Prefer not to answer'' and ``I'm not
sure'' were retained as substantive categories in all frequency tables, excluded from ordinal numeric tests
because they carry no rank position, and treated as item-missing rather than scored as zero in composite
indices. Subgroup cells containing fewer than five respondents are suppressed and are not back-calculated;
counts of exactly five are reported exactly, a disclosure decision the research team made deliberately on
the grounds that no respondent is identifiable from these aggregates.

Because the platform populated no duplicate-detection metadata and collected no IP or device information,
duplicate submissions were screened manually on the respondent-chosen nickname and password. One pair
shared both fields, and its earlier member was already excluded as an abandoned submission, so no analytic
row required exclusion on duplicate grounds; repeated nicknames or passwords without a matching second
field involved common Spanish words and were judged coincidental (Supplementary Material). Two respondents
failed at least one attention check and were flagged rather than dropped, per protocol.

\subsection{Statistical approach}

The outcome distribution was strongly ceiling-weighted and several exposures had thin or unbalanced
categories, so all tests were nonparametric: Mann--Whitney $U$ for two-group ordinal comparisons,
Kruskal--Wallis $H$ for three or more groups, permutation $\chi^2$ with 10,000 label permutations (seed
20260715) for thin contingency tables, and Fisher's exact test for $2\times2$ tables at very small $n$.
Effect sizes were rank-biserial $r_{rb}$ with the common-language effect size, epsilon-squared
$\varepsilon^2$, and Cram\'er's $V$ respectively, with Wilson 95\% confidence intervals for prevalences.
The bias-corrected $\varepsilon^2$ formula can return small negative values when an effect is
indistinguishable from zero; these are displayed truncated to 0.00 (marked $\dagger$), applied uniformly.

The preparation-to-satisfaction comparison was pre-declared and is reported uncorrected, with a tie-aware
permutation test as a robustness check because the outcome is heavily tied. All other bivariate tests were
assigned to four prespecified secondary and exploratory families---bivariate, substance-stratified,
safety-domain, and legal-knowledge---corrected jointly within family by the Benjamini--Hochberg
false-discovery-rate procedure and kept separate because they share neither exposures nor outcome domains.
One multivariable proportional-odds model was fitted as an exploratory supplement. Analyses were conducted
in Python; the full pipeline is available (see Data availability).

\subsection{Protocol deviations}

Three deviations are reported transparently. First, the original protocol specified a four-arm preparation
$\times$ integration design. This is not estimable from the collected data: the instrument contains no
clean participant integration-received item, and the preparation split is badly unbalanced. Preparation was
therefore examined against the outcome directly, and no integration measure is reconstructed anywhere in
this paper.

Second, the protocol specified crossing facilitator training and protocol items against participant-perceived
safety. Because those items are facilitator self-report from a different and unlinked set of respondents,
crossing them would have required pairing unrelated people. Two participant-answerable proxies were
substituted---prior knowledge of the facilitator's background and whether screening questions were
asked---and facilitator self-report is presented separately and descriptively.

Third, the protocol specified a composite harm-reduction score mixing facilitator-side and participant-side
items, which for the same reason could not validly be built; two branch-specific composites were
constructed instead. Two skip-logic errors in the instrument were also identified and are documented in the
Supplementary Material.

\section{Results}

\subsection{Respondents}

The analytic sample of 73 adults skewed middle-aged and highly educated (Table~\ref{tab:demo}). Among the
72 respondents with a valid age, the mean was 42.4 years and the median 41 (SD 11.9, range 21--71, IQR
34--50). Of the 71 answering the gender item, 35 (49\%) identified as women and 30 (42\%) as men, with
non-binary and prefer-not-to-answer cells suppressed. Of the 72 answering the education item, 34 (47\%)
reported a bachelor's degree and 19 (26\%) a master's, and of the 72 answering the economic item, 44
(61\%) described their situation as comfortable against 6 (8\%) as difficult. Sixty-eight respondents
(94\% of the 72 answering) reported residing in Puerto Rico; residency was treated as a descriptive
covariate rather than an inclusion filter, and the presence of respondents living elsewhere tempers any
reading of the sample as describing Puerto Rico specifically.

Free-text engagement was substantial relative to sample size: 37 open-text fields yielded 285 non-empty
answers, roughly three-quarters of them in Spanish, most often on an experience highlight (58 answers),
opinions on safety (54), and preparation guidance received (45). A coding scaffold for the research team's
thematic analysis accompanies these data (Supplementary Material); the corpus is not interpreted here.

\subsection{The ceremonial context}

Participants reported extensive but not recent ceremonial histories. Among the 67 answering, the median
number of lifetime ceremonies attended was 6 (mean 15.1, SD 35.7, range 1--250, IQR 3--12), a long right
tail that makes the median the appropriate summary, while ceremonies in the past 12 months were far fewer
in the same 67 (median 1, mean 1.3, SD 1.7, range 0--8, IQR 0--2). Of the 66 answering the recency item, 30
(45\%) reported their last ceremony more than 12 months earlier.

Across all ceremonies attended, the substances most often reported by the 67 participants answering were
magic mushrooms (50, 75\%), ayahuasca (43, 64\%), tobacco (38, 57\%), and cacao (23, 34\%), with DMT, LSD,
MDMA, and ketamine less frequent (Figure~\ref{fig:subs}). For the most recent ceremony specifically,
individually reportable groups were magic mushrooms ($n=31$) and ayahuasca ($n=23$), with every remaining
substance individually below the suppression floor and pooled into a combined group ($n=11$) under a
prespecified rule.

The most recent ceremony was typically a discrete, multi-hour, outdoor event (Table~\ref{tab:context}). Of
the 65 answering the format item, 48 (74\%) reported a single ceremony rather than a retreat or series; of
the 67 answering the place item, 37 (55\%) reported a natural setting such as a forest, beach, or mountain;
and of the 65 answering the duration item, 19 (29\%) reported 4--6 hours and a further 19 (29\%) reported
an overnight ceremony of roughly 12 hours or more.

\subsection{Preparation, screening, and prior knowledge}

Preparation was near-universal among those answering (Table~\ref{tab:safety}): of the 67 answering, 55
(82\%) reported preparing and 6 (9\%) preparing sometimes, against 6 (9\%) reporting no self-directed
preparation. This imbalance is why the four-arm design was abandoned and it constrains every preparation
analysis below. Participant-observable facilitator practices were common but not uniform. Of the 67
answering, 54 (81\%) reported knowing about the facilitator's background or training beforehand and 8
(12\%) knowing a little, with the remaining cells suppressed. On screening, of the 66 answering, 39 (59\%)
reported that the facilitator asked screening questions beforehand, 15 (23\%) that none were asked, and 7
(11\%) were unsure. Disclosure outside the ceremonial context was uncommon: of the 65 answering, 8 (12\%)
reported having spoken to a doctor, therapist, or other health professional about their participation
beforehand, against 54 (83\%) who had not.

A participant-side harm-reduction practice index counting these three practices (preparation, prior
knowledge of the facilitator's background, screening received; range 0--3) was constructed for the 54
participants with complete data on all three. Thirty-five (65\%) reported all three and 16 (30\%) reported
two; \emph{none} reported zero. Internal consistency was low ($\alpha = 0.14$, 95\% CI $-0.36$--0.47),
indicating that the items capture distinct protective practices rather than one latent construct, so the
score is read as a count of practices present rather than as a validated scale.

\subsection{Experience and perceived safety}

Reported satisfaction was overwhelmingly positive. Of the 72 participants, 63 answered the satisfaction
item, and the distribution was strongly ceiling-weighted with a median of 6 and the mode at the scale
maximum (Figure~\ref{fig:good}). Perceived in-ceremony comfort and safety was similarly high and similarly
undifferentiated, returning a median of 3 in every subgroup tested and differing neither by prior knowledge
of the facilitator's background ($H=2.24$, $p=0.326$, $\varepsilon^2=0.00$) nor by whether screening
questions were asked ($H=2.40$, $p=0.301$, $\varepsilon^2=0.01$). Trust in the facilitator was the one
participant-side comparison differing at the uncorrected level: median 4 (IQR 3--4) among the 54 who knew
the facilitator's background beforehand against 3 (IQR 2--3) among the 8 who knew a little, with a
suppressed third group ($H=8.15$, $p=0.017$, $\varepsilon^2=0.10$); it did not survive correction
($q=0.102$). Satisfaction did not differ significantly by most-recent-ceremony substance (ayahuasca median
5, mushrooms median 6, combined group median 6; $H=4.98$, $p=0.083$, $q=0.166$).

The safety domain produced the sample's most consequential descriptive finding. Of the 67 participants
answering the non-consensual contact item, 61 (91\%) reported no such contact and 5 (7\%) reported that it
occurred, with a small uncertain group suppressed; the within-sample proportion was 7.5\% (Wilson 95\% CI
3.2--16.3). This describes these 67 respondents and nothing beyond them. Non-consensual contact was
associated with prior knowledge of the facilitator's background at the uncorrected level (permutation
$\chi^2 = 18.90$, $p=0.047$, $V=0.38$, $n=67$) but not with screening (permutation $\chi^2 = 13.64$,
$p=0.094$, $V=0.32$, $n=66$); neither survived correction (both $q=0.236$), nor did the comparison by
substance ($p=0.486$, $q=0.486$) or by the harm-reduction index ($p=0.641$, $q=0.641$). The research team
elected to publish these cross-tabulations under the same suppression convention applied to every other
thin-cell table rather than withhold them: withholding would remove the only evidence this study can offer
on whether such reports cluster with any measured practice, and the masking convention leaves no cell small
enough to plausibly identify a respondent. Every result in this domain is hypothesis-generating and is
powered neither to confirm nor to rule out an association.

One measurement observation belongs here: an item on perceived positive life impact had zero variance, with
all 67 respondents answering selecting the same option.

\subsection{Facilitator practice}

Eight respondents identified as facilitators. This is a different, unlinked set from the participants
above, nearly every category cell falls at or below the suppression floor, and no group comparison at this
base is meaningful; the material is descriptive context only. Preparation for the role was informal: 6 of
the 8 (75\%) reported personal experience and self-study, with mentorship and family or ancestral tradition
each reported by suppressed numbers, and \emph{no} facilitator reported formal certification or academic
coursework. Six of 8 (75\%) reported an initial conversation or interview with prospective participants,
and 5 of 8 (62\%) reported using singing, music, or sound to guide someone through a difficult experience.
On emergency planning, all responses at a base of 7 fell into informal or conditional categories, with no
facilitator selecting the formal option. Seven facilitators answered the crisis-signs item; individual
sub-items, including suicidal ideation and panic attack, are each suppressed and establish only that such
events were witnessed and fall within scope. Association tests within this branch returned undefined odds
ratios at $n=7$ (Supplementary Material).

\subsection{Preparation and reported satisfaction}

Among the 63 participants answering the satisfaction item, 58 reported some preparation (median 6, IQR
5--6) and 5 reported none (median 4, IQR 3--5; Figure~\ref{fig:prep}); one further respondent reporting no
preparation did not answer the satisfaction item, which is why this base is 5 rather than the 6 in
Table~\ref{tab:safety}. The pre-declared Mann--Whitney
comparison returned $U=218.0$, $p=0.041$, $r_{rb}=0.50$, and a common-language effect size of 0.75, with
preparation associated with higher reported satisfaction; a tie-aware permutation test gave $p=0.024$. The
structural weakness of this comparison, however, precedes its result: the no-preparation group contains
five respondents. The prespecified three-level sensitivity analysis was non-significant ($H=4.24$,
$p=0.120$, $\varepsilon^2=0.04$; $q=0.359$), the practice index was unassociated with satisfaction
($H=2.09$, $p=0.352$; $q=0.422$), and an exploratory proportional-odds model adjusting preparation for age
and facilitator trust placed preparation short of conventional significance (OR 5.21, 95\% CI 0.84--32.41,
$p=0.077$) while trust reached it (OR 2.20, 95\% CI 1.06--4.53, $p=0.033$) at roughly 12 cases per
estimated parameter (Supplementary Material). The comparison is suggestive and exploratory, not
confirmatory. No test in any corrected family was significant at $q<0.05$.

\section{Discussion}

\subsection{A structurally informal practice landscape}

The most stable picture this sample supports is a practice organized around informal, relationally
transmitted expertise rather than credentialed or documented protocol. Among the eight facilitators who
self-reported, training came from personal experience, self-study, mentorship, and family tradition; none
reported formal certification or academic coursework, and emergency planning was uniformly informal or
conditional. This sits at some distance from the structured preparation, screening, and difficult-experience
protocols that harm-reduction frameworks for psychedelic settings assume
\citep{gorman2021psychedelic,pilecki2021ethical,dutton2025harm}, and it is consistent with survey evidence
that people who use psychedelics support clearer practice standards \citep{kruger2023preferences}.

On the participant side the informality reads differently. Preparation was near-universal among those
answering, but roughly a quarter of participants answering the screening item reported that no screening
questions were asked and a further group was uncertain. The asymmetry is worth separating: preparation is
something participants largely do, whether self-directed or instructed, whereas screening must be done
\emph{to} them, and was substantially less consistent in what these respondents reported. If that
asymmetry holds in a larger sample it would locate the intervention target on the facilitator side---but
this sample cannot establish whether it holds, and eight facilitators cannot characterize screening
practice.

The demographic profile---middle-aged, university-educated, economically comfortable---resembles that of
ceremony-attender cohorts described elsewhere \citep{teixeira2026ayahuasca,pagni2025longterm}, of large
international psychedelic-survey samples \citep{kopra2023investigation,robinson2026field}, and of the
psilocybin-specific Puerto Rico survey conducted by this team \citep{velezrodriguez2026psilocybin}. That
convergence is best read as evidence about who responds to online psychedelic surveys, not about who
attends ceremonies.

\subsection{A ceiling that limits what satisfaction can reveal}

Reported satisfaction clustered at the top of its scale, and a perceived-life-impact item had no variance
at all. This is a substantive finding---these respondents overwhelmingly described their most recent
ceremony positively, consistent with naturalistic cohorts reporting predominantly positive appraisals
\citep{nayak2023naturalistic,ruffell2021ceremonial}---but it also constrains the analysis severely. An
outcome with almost no variance in the lower range cannot discriminate between groups, and the pattern of
null results across the corrected families is consistent with that constraint rather than with the absence
of any association. Ceiling behavior is a known property of self-report scales for hallucinogen experience
rather than a peculiarity of this instrument \citep{bouso2016measuring}. The design implication is that
satisfaction, measured this way, is close to exhausted as an outcome in this population; future work needs
measures with headroom in both directions before group comparisons on experience quality can be
informative.

\subsection{Safety signals the design can identify but not quantify}

Five of 67 participants answering reported non-consensual physical contact, and facilitators reported
having witnessed crisis signs, though every individual sub-item was suppressed and no rate was estimable at
a base of seven. These two findings come from different, unlinked reporters and cannot be combined. That
such reports appear at all is consistent with the only survey-based evidence on boundary violations in
naturalistic psychedelic settings \citep{kruger2025psychedelic} and with the community harm-reduction
literature that has developed in response \citep{peluso2020reflections}. The item used here was
deliberately broad, covering unexpected touch and uncertain or unremembered consent, so the figure should
not be read against a narrower legal or clinical standard.

The study establishes that these events occur among respondents and that an instrument of this kind can
capture them. It cannot establish how often they occur, in whom, or under what conditions. The
cross-tabulations against participant-side proxies are powered neither to confirm nor to rule out an
association, and the one comparison reaching uncorrected significance did not survive correction; treating
any of them as an actionable signal would overstate them considerably. The safety domain justified its
inclusion by identifying a measurement priority, not by producing an estimate.

\subsection{The preparation finding as a design lesson}

The preparation-to-satisfaction comparison is this study's pre-declared primary analysis and its most
fragile result, and those two facts should be held together rather than traded off. The direction was
positive and the effect sizes were not small, but the no-preparation group contained five respondents
answering the outcome, the three-level sensitivity analysis was non-significant, and the adjusted model
placed preparation short of significance with a confidence interval spanning nearly two orders of
magnitude. Nothing here establishes the finding.

The more durable contribution is diagnostic. The preparation split was unbalanced enough to dissolve the
original four-arm design and it remains the binding constraint on the primary analysis. Any successor study
intending to examine preparation will need either a sampling strategy that reaches unprepared participants
in usable numbers or a question that does not depend on that contrast. Reporting the deviation rather than
absorbing it quietly is what makes the constraint visible to the next design.

\subsection{Limitations}

The sample was non-probability, unweighted, and self-selected into an online survey; the concentration in
higher education categories and favorable economic self-ratings describe who responded, not who
participates. Nothing here supports a statement about how common ceremonial entheogen use is in Puerto
Rico or about people who use entheogens in Puerto Rico generally. Roughly 6\% of those answering the
residency item lived outside Puerto Rico, further tempering any Puerto Rico--specific reading.

Role subsamples were small and unevenly sized, and the facilitator branch of eight supports no rate and no
association test. Participant and facilitator respondents are unlinked, which is not a minor analytic
inconvenience: it eliminated the protocol's integration-proxy comparison, its perceived-safety-by-training
analysis, and its mixed harm-reduction composite, and it means no statement in this paper connects what
facilitators reported doing to what participants reported experiencing.

All measures were single-timepoint self-report, and for most respondents the most recent ceremony was
distant---30 of 66 answering reported more than 12 months---so recall error is plausible and unmeasurable
here. The instrument contains no integration measure, is unvalidated, and produced mixed first reliability
evidence alongside two documented skip-logic errors (Supplementary Material).

\subsection{Conclusion}

Within a self-selected sample of 73 adults, ceremonial entheogen use in Puerto Rico was described by a
middle-aged, university-educated, overwhelmingly Puerto Rico--resident group whose most recent ceremonies
were typically single, multi-hour, outdoor events involving psilocybin mushrooms or ayahuasca, and who
described those ceremonies in strongly positive terms. Preparation was near-universal while pre-ceremony
screening was not, and the small facilitator subsample described training and emergency planning as
informal rather than credentialed. Five of 67 participants answering reported non-consensual physical
contact, and facilitators reported witnessing crisis events at a base too small for any rate to be
estimated. These findings describe this sample alone. They support no estimate of prevalence and no
inference about what any practice produces, and their principal value lies in identifying which measures,
sampling strategies, and instrument corrections a properly powered successor study would require.

\section*{Author contributions}
% UNCONFIRMED — drafted from the roles described in the study protocol. Requires author sign-off.
YOOO conceived the study, led protocol development and ethics submission, and supervised the project.
YOOO, JCVR, JMHT, JMT, PRF, CR, ANGR, and AIRM contributed to instrument development and review. JCVR led
data curation, formal analysis, and manuscript drafting. JMHT contributed to methodology and analysis
review. JMT, PRF, and CR contributed to interpretation and clinical/community context. ANGR and AIRM
contributed to recruitment and data collection. All authors reviewed and approved the final manuscript.

\section*{Funding}
% UNCONFIRMED — no funding source is named in any study document. Requires author sign-off.
This research received no specific grant from any funding agency in the public, commercial, or
not-for-profit sectors.

\section*{Competing interests}
% UNCONFIRMED — requires disclosure from each author. Note that several authors are affiliated with
% psychedelic-focused organizations, which may be disclosable under JPD policy.
The authors declare no competing interests.

\section*{Data availability}
% UNCONFIRMED — confirm whether de-identified data can be shared, given the protocol's retention terms.
The analysis code that generates every result reported here is available from the corresponding author.
De-identified survey data are available from the corresponding author on reasonable request, subject to the
terms of the approved protocol.

\section*{Generative AI use declaration}
% UNCONFIRMED — verify scope before submission.
Large language model tooling (Anthropic Claude) was used to assist with statistical analysis code
development and with drafting and editing of this manuscript. All analyses were specified by the authors,
all outputs were verified against the source data by the authors, and the authors take full responsibility
for the content of this article.

\bibliography{../references/references}

\clearpage

\section*{Tables and Figures}

\begin{table}[htbp]
\centering
\caption{Sociodemographic characteristics of the analytic sample ($N=73$). Bases differ by item; cells with
$n<5$ are suppressed and are not back-calculated. Percentages are of respondents answering.}
\label{tab:demo}
\setstretch{1.1}
\begin{tabular}{p{3.6cm}p{6.4cm}rr}
\toprule
Characteristic & Category & $n$ & \% answering \\
\midrule
Age, years ($n=72$) & Mean 42.4 (SD 11.9); median 41 (IQR 34--50); range 21--71 & & \\
\addlinespace
Gender ($n=71$)     & Woman & 35 & 49 \\
                    & Man & 30 & 42 \\
                    & Non-binary & $<5$ & 6 \\
                    & Prefer not to answer & $<5$ & 3 \\
\addlinespace
Education ($n=72$)  & Bachelor's degree & 34 & 47 \\
                    & Master's degree & 19 & 26 \\
                    & Associate's degree & 7 & 10 \\
                    & High school diploma / GED & 6 & 8 \\
                    & Doctorate & $<5$ & 4 \\
                    & Prefer not to answer & $<5$ & 4 \\
\addlinespace
Economic situation ($n=72$) & Comfortable & 44 & 61 \\
                    & Tight & 20 & 28 \\
                    & Difficult & 6 & 8 \\
                    & Prefer not to answer & $<5$ & 3 \\
\addlinespace
Country of residence ($n=72$) & Puerto Rico & 68 & 94 \\
                    & United States & $<5$ & 4 \\
                    & Spain & $<5$ & 1 \\
\bottomrule
\end{tabular}
\end{table}

\begin{table}[htbp]
\centering
\caption{Most-recent-ceremony context, participant-reported. Participant branch only (base 72); bases
differ by item because structural skips are reported separately from item nonresponse. Cells with $n<5$ are
suppressed.}
\label{tab:context}
\setstretch{1.1}
\begin{tabular}{p{4.2cm}p{5.8cm}rr}
\toprule
Item & Category & $n$ & \% answering \\
\midrule
Time since last ceremony ($n=66$) & More than 12 months ago & 30 & 45 \\
                                  & 2 to 6 months ago & 20 & 30 \\
                                  & 7 to 12 months ago & 11 & 17 \\
                                  & During the past month & $<5$ & 6 \\
\addlinespace
Place ($n=67$)   & Natural setting (forest, beach, mountain) & 37 & 55 \\
                 & Private home & 15 & 22 \\
                 & Retreat center or wellness facility & 9 & 13 \\
\addlinespace
Format ($n=65$)  & Single ceremony & 48 & 74 \\
                 & Retreat or series of ceremonies & 13 & 20 \\
\addlinespace
Duration ($n=65$) & 4 to 6 hours & 19 & 29 \\
                 & Overnight ($\sim$12+ hours) & 19 & 29 \\
                 & 7 to 12 hours & 14 & 22 \\
                 & 1 to 3 hours & 7 & 11 \\
\bottomrule
\end{tabular}
\end{table}

\begin{table}[htbp]
\centering
\caption{Participant-reported protective practices and safety outcomes. Bases are per item and differ
because each is asked under its own skip condition. Cells with $n<5$ are suppressed; counts of exactly 5
are reported exactly.}
\label{tab:safety}
\setstretch{1.1}
\begin{tabular}{p{5.2cm}p{4.8cm}rr}
\toprule
Item & Category & $n$ & \% answering \\
\midrule
Self-directed preparation ($n=67$) & Yes & 55 & 82 \\
                & Sometimes & 6 & 9 \\
                & No & 6 & 9 \\
\addlinespace
Knew facilitator's background ($n=67$) & Yes & 54 & 81 \\
                & A little & 8 & 12 \\
                & No & $<5$ & 6 \\
                & Prefer not to answer & $<5$ & 1 \\
\addlinespace
Facilitator asked screening questions ($n=66$) & Yes & 39 & 59 \\
                & No & 15 & 23 \\
                & Not sure & 7 & 11 \\
                & Prefer not to answer & 5 & 8 \\
\addlinespace
Spoke to a health professional beforehand ($n=65$) & No & 54 & 83 \\
                & Yes & 8 & 12 \\
                & Prefer not to answer / not sure & $<5$ & 5 \\
\addlinespace
Non-consensual physical contact ($n=67$) & No & 61 & 91 \\
                & Yes & 5 & 7 \\
                & Not sure & $<5$ & 1 \\
\addlinespace
Harm-reduction practice index ($n=54$) & All three practices & 35 & 65 \\
                & Two practices & 16 & 30 \\
                & One practice & $<5$ & 6 \\
                & No practices & 0 & 0 \\
\bottomrule
\end{tabular}
\end{table}

\begin{figure}[htbp]
\centering
\fbox{\parbox[c][6cm][c]{0.85\textwidth}{\centering\itshape
Figure to be generated in Phase 4 (sample-construction flow): 100 submissions $\rightarrow$ 13 abandoned
excluded $\rightarrow$ 87 reaching the role fork $\rightarrow$ 14 screen-outs excluded $\rightarrow$ $N=73$
analytic sample $\rightarrow$ role fork into 65 participant-only, 7 dual-role, 1 facilitator-only, with the
unlinked-branch structure marked.}}
\caption{Sample construction and role structure. The participant and facilitator branches are unlinked: no
identifier connects a participant to the facilitator who ran their ceremony.}
\label{fig:flow}
\end{figure}

\begin{figure}[htbp]
\centering
\includegraphics[width=0.85\textwidth]{../figures/fig2_substances.png}
\caption{Substances used across all ceremonies attended, participant-reported (select-multiple; base $n=67$
participants answering; percentages need not sum to 100\%). Distinct from the facilitator-led substance
variable and from the most-recent-ceremony substance variable.}
\label{fig:subs}
\end{figure}

\begin{figure}[htbp]
\centering
\includegraphics[width=0.85\textwidth]{../figures/fig3_ceremony_good.png}
\caption{Distribution of reported satisfaction with the most recent ceremony (0 = very unpleasant to 6 =
very pleasant). Base: 63 of 72 participants answering; the remainder is item nonresponse. Median 6.}
\label{fig:good}
\end{figure}

\begin{figure}[htbp]
\centering
\includegraphics[width=0.85\textwidth]{../figures/fig4_prep_outcome.png}
\caption{Reported satisfaction by self-directed preparation status. Any preparation, $n=58$ answering the
outcome (median 6, IQR 5--6); no preparation, $n=5$ (median 4, IQR 3--5). \textbf{Exploratory:} the
no-preparation group contains five respondents and this comparison is suggestive, not confirmatory.}
\label{fig:prep}
\end{figure}

\end{document}
